\section{Tổng quan tài liệu}
%Ghi rõ và ưu nhược điểm
\subsection{Các phương pháp Học máy hiện có để giải quyết bài toán}
%Ghi rõ và ưu nhược điểm
Dữ liệu mà nhóm nghiên cứu có các tính chất đặc thù như quy mô dữ liệu khổng lồ, tính chất chuỗi thời gian, tồn tại nhiễu cảm biến diện rộng và sự mất cân bằng giữa 2 lớp nghiêm trọng.

\subsubsection{Nhóm Phân Lớp Có Giám Sát (Supervised Classification)}
Do dữ liệu đã có gán nhãn cho biến \texttt{anomaly}( $0$ biểu thị trạng thái bình thường, $1$ biểu thị trạng thái bất thường), họ thuật toán dựa trên cây quyết định (Tree-based Models) và phương pháp học tập hợp (Ensemble Learning) thường thể hiện hiệu suất vượt trội.

\paragraph{XGBoost (eXtreme Gradient Boosting)}\hspace{1cm}

XGBoost tối ưu hóa hàm mục tiêu bằng phương pháp khai triển Taylor bậc hai, cho phép hội tụ nhanh hơn. Cấu trúc cây được phát triển theo từng tầng (Level-wise growth). Mô hình cuối cùng là tổng hợp tuần tự của $K$ cây quyết định:
\begin{equation*}
    \hat{y}_i = \sum_{k=1}^{K} f_k(\mathbf{x}_i), \quad f_k \in \mathcal{F}
\end{equation*}
trong đó $\mathcal{F}$ là không gian của các cây hồi quy CART. Tại mỗi bước lặp $t$, hàm mục tiêu được tối ưu bằng khai triển Taylor bậc hai:
\begin{equation*}
    \mathcal{L}^{(t)} \simeq \sum_{i=1}^{n} \left[ g_i f_t(\mathbf{x}_i) + \frac{1}{2} h_i f_t^2(\mathbf{x}_i) \right] + \Omega(f_t)
\end{equation*}
với $g_i = \partial_{\hat{y}^{(t-1)}} \ell(y_i, \hat{y}^{(t-1)})$ là gradient bậc nhất, $h_i = \partial^2_{\hat{y}^{(t-1)}} \ell(y_i, \hat{y}^{(t-1)})$ là gradient bậc hai (Hessian) của hàm mất mát và $\Omega(f_t)$ là hàm phạt cấu trúc cây:
\begin{equation*}
    \Omega(f) = \gamma T + \frac{1}{2} \lambda \sum_{j=1}^{T} w_j^2 + \alpha \sum_{j=1}^{T} |w_j|
\end{equation*}
trong đó $T$ là số lá, $w_j$ là trọng số tại lá $j$, $\gamma$ phạt số lượng lá, $\lambda$ là hệ số điều chuẩn $L_2$ và $\alpha$ là hệ số điều chuẩn $L_1$.\\

\textbf{Ưu điểm:}
\begin{itemize}
    \item \textbf{Cơ chế điều chuẩn (Regularization):} Tích hợp các thành phần phạt $L_1$ ($\alpha$) và $L_2$ ($\lambda$) trực tiếp vào hàm mục tiêu $\Omega(f)$, giúp hạn chế hiệu quả hiện tượng overfitting khi mô hình tiếp xúc với dữ liệu nhiễu từ các cảm biến đo lường.
    \item \textbf{Xử lý mất cân bằng lớp:} Hỗ trợ điều chỉnh trọng số thông qua siêu tham số \texttt{scale\_pos\_weight}, gán penalty cao hơn cho các trường hợp phân loại sai lớp thiểu số.
    \item \textbf{Khai triển Taylor bậc hai:} Sử dụng cả gradient bậc nhất $g_i$ và Hessian $h_i$ giúp xấp xỉ hàm mất mát chính xác hơn, tạo ra các bước cập nhật tối ưu hơn so với các phương pháp Gradient Boosting truyền thống chỉ dùng gradient bậc nhất.
    \item \textbf{Xử lý giá trị thiếu tự động:} Thuật toán tự học hướng mặc định (default direction) tại mỗi nút chia cho các giá trị bị khuyết, phù hợp với dữ liệu cảm biến thực tế thường có nhiều giá trị bị thiếu.
\end{itemize}

\textbf{Nhược điểm:}
\begin{itemize}
    \item \textbf{Độ phức tạp không gian và thời gian:} Thuật toán phân nhánh tham lam chính xác (Exact Greedy Algorithm) đòi hỏi duyệt qua toàn bộ dữ liệu liên tục. Mặc dù đã có các cải tiến như Histogram-based, XGBoost vẫn yêu cầu lượng lớn RAM và thời gian tính toán khi xử lý tập dữ liệu có quy mô lớn.
    \item \textbf{Phát triển cây Level-wise:} Chiến lược phát triển theo tầng buộc tất cả các nút ở cùng một độ sâu phải được mở rộng, dẫn đến việc phân chia tại nhiều nút không cần thiết, gây lãng phí tài nguyên tính toán.
\end{itemize}

\paragraph{LightGBM (Light Gradient Boosting Machine)}\hspace{1cm}

LightGBM được thiết kế chuyên biệt để tối ưu hóa hiệu suất trên dữ liệu quy mô lớn thông qua hai kỹ thuật cốt lõi: GOSS (Gradient-based One-Side Sampling) và EFB (Exclusive Feature Bundling). Cây được phát triển theo cơ chế ưu tiên lá (Leaf-wise growth).

Kỹ thuật GOSS ước lượng độ lợi thông tin (Information Gain) bằng cách giữ lại toàn bộ các mẫu có gradient lớn (top $a\%$) và lấy mẫu ngẫu nhiên từ các mẫu có gradient nhỏ (tỷ lệ $b\%$), sau đó nhân trọng số bù:
\begin{equation*}
    \widetilde{V}_j(d) = \frac{1}{n} \left( \frac{\left(\sum_{x_i \in A_l} g_i + \frac{1-a}{b} \sum_{x_i \in B_l} g_i\right)^2}{n_l^j(d)} + \frac{\left(\sum_{x_i \in A_r} g_i + \frac{1-a}{b} \sum_{x_i \in B_r} g_i\right)^2}{n_r^j(d)} \right)
\end{equation*}
trong đó $A$ là tập các mẫu có gradient lớn, $B$ là tập mẫu ngẫu nhiên từ phần gradient nhỏ, $l$ và $r$ ký hiệu nhánh trái và phải tại ngưỡng $d$ của đặc trưng $j$.\\

Kỹ thuật EFB (Exclusive Feature Bundling) gộp các đặc trưng thưa có tính loại trừ lẫn nhau thành các gói (bundles), giảm số lượng đặc trưng hiệu dụng từ $p$ xuống số lượng gói $p' \ll p$.

\textbf{Ưu điểm:}
\begin{itemize}
    \item \textbf{Hiệu suất tính toán ưu việt:} Kỹ thuật Histogram giúp giảm độ phức tạp thời gian tìm kiếm điểm chia nhánh xuống $O(\text{số bin} \times \text{số đặc trưng})$, mang lại tốc độ huấn luyện nhanh gấp 3 đến 10 lần so với XGBoost chuẩn.
    \item \textbf{Xử lý biến định danh trực tiếp:} Khai thác hiệu quả các đặc trưng phân loại có lực lượng cao (high-cardinality) như \texttt{building\_id} hay \texttt{meter} bằng cách phân chia tối ưu trên các danh mục mà không cần mã hóa One-Hot, giúp kiểm soát chiều dữ liệu.
    \item \textbf{Tiết kiệm bộ nhớ:} Kỹ thuật EFB giảm đáng kể số chiều hiệu dụng, trong khi biểu diễn Histogram chỉ lưu trữ các bin rời rạc thay vì toàn bộ giá trị liên tục, giúp tiết kiệm RAM khi xử lý dữ liệu hàng chục triệu dòng.
    \item \textbf{GOSS giữ độ chính xác khi lấy mẫu:} Thay vì lấy mẫu đồng đều, GOSS ưu tiên các mẫu có gradient lớn (đóng góp nhiều nhất vào việc học), nên giảm dữ liệu mà vẫn duy trì chất lượng mô hình.
\end{itemize}

\textbf{Nhược điểm:}
\begin{itemize}
    \item \textbf{Rủi ro quá khớp cục bộ:} Sự phát triển theo chiều sâu lá (Leaf-wise) có thể dẫn đến việc mô hình học các mẫu nhiễu cục bộ. Cần giới hạn nghiêm ngặt các tham số \texttt{max\_depth} và \texttt{min\_data\_in\_leaf}.
    \item \textbf{Nhạy cảm với dữ liệu ít:} Chiến lược Leaf-wise tạo ra các cây sâu không cân bằng, dễ overfitting trên các tập dữ liệu nhỏ. Tuy nhiên, hạn chế này ít ảnh hưởng trong bài toán hiện tại do quy mô dữ liệu rất lớn.
\end{itemize}

\paragraph{CatBoost (Categorical Boosting)}\hspace{1cm}

Thuật toán xây dựng các cây quyết định đối xứng (Oblivious Trees) và sử dụng cơ chế Ordered Target Encoding. Đặc điểm cốt lõi của CatBoost là sử dụng hoán vị ngẫu nhiên (random permutation) $\sigma$ để tính giá trị mã hóa cho biến định danh, tránh rò rỉ thông tin mục tiêu:
\begin{equation*}
    \hat{x}_{\sigma_p}^i = \frac{\sum_{j=1}^{p-1} [x_{\sigma_j}^i = x_{\sigma_p}^i] \cdot y_{\sigma_j} + a \cdot P}{\sum_{j=1}^{p-1} [x_{\sigma_j}^i = x_{\sigma_p}^i] + a}
\end{equation*}
trong đó $\sigma_p$ là chỉ số của mẫu thứ $p$ trong hoán vị, $x^i$ là giá trị đặc trưng định danh thứ $i$, $y$ là nhãn mục tiêu, $P$ là prior (thường là trung bình toàn cục của $y$) và $a > 0$ là tham số Laplace smoothing.

Cây đối xứng (Oblivious Tree) sử dụng cùng một điều kiện chia tại tất cả các nút ở cùng một độ sâu, tạo ra cấu trúc cây cân bằng hoàn hảo với $2^d$ lá (với $d$ là chiều sâu). Dự đoán trở thành phép tra bảng $O(d)$.\\

\textbf{Ưu điểm:}
\begin{itemize}
    \item \textbf{Ngăn chặn rò rỉ dữ liệu (Target Leakage):} Ordered Target Encoding tính toán giá trị trung bình mục tiêu dựa trên các mẫu quan sát trong quá khứ nhân tạo theo hoán vị $\sigma$, giúp mã hóa các biến định danh một cách an toàn và mạnh mẽ mà không gây rò rỉ thông tin nhãn.
    \item \textbf{Hiệu năng suy luận (Inference Time):} Cấu trúc cây đối xứng cho phép mô hình được biểu diễn dưới dạng các bảng tra cứu với độ phức tạp suy luận $O(d)$, tăng tốc độ dự đoán lên nhiều lần so với các cây không đối xứng.
    \item \textbf{Ordered Boosting:} Sử dụng nhiều hoán vị khác nhau trong quá trình huấn luyện, giảm thiên kiến (bias) do rò rỉ thông tin trong Gradient Boosting truyền thống, đặc biệt hiệu quả trên dữ liệu có nhiều biến phân loại như \texttt{building\_id}, \texttt{primary\_use}.
    \item \textbf{Kháng overfitting tự nhiên:} Cấu trúc cây đối xứng có ít tham số hơn cây tự do (chỉ $d$ điều kiện chia thay vì $2^d - 1$ nút nội), đóng vai trò như một dạng điều chuẩn ngầm.
\end{itemize}

\textbf{Nhược điểm:}
\begin{itemize}
    \item \textbf{Thời gian huấn luyện trên dữ liệu số thực:} Quy trình mã hóa Ordered Target Encoding và xây dựng cây đối xứng khiến CatBoost mất nhiều thời gian huấn luyện hơn so với LightGBM, đặc biệt khi bộ dữ liệu chứa lượng lớn các biến liên tục.
    \item \textbf{Giới hạn biểu diễn của cây đối xứng:} Ràng buộc cùng một điều kiện chia tại mỗi tầng có thể hạn chế khả năng biểu diễn các mẫu phức tạp, đòi hỏi cây phải sâu hơn để đạt cùng mức chính xác so với cây tự do.
\end{itemize}

\paragraph{Random Forest}\hspace{1cm}

Random Forest là thuật toán học tập hợp Bagging (Bootstrap Aggregating), xây dựng $B$ cây quyết định độc lập trên các tập dữ liệu bootstrap. Mỗi cây được huấn luyện trên tập $\mathcal{D}_b$ được lấy mẫu ngẫu nhiên có hoàn lại (with replacement) từ tập gốc $\mathcal{D}$, với $|\mathcal{D}_b| = |\mathcal{D}| = n$.

Tại mỗi nút phân chia, thay vì xét toàn bộ $p$ đặc trưng, Random Forest chỉ chọn ngẫu nhiên một tập con $m$ đặc trưng và tìm điểm chia tối ưu theo tiêu chí Gini Impurity:
\begin{equation*}
    \text{Gini}(S) = 1 - \sum_{k=1}^{C} p_k^2
\end{equation*}
trong đó $C$ là số lớp và $p_k$ là tỷ lệ mẫu thuộc lớp $k$ trong tập $S$. Số đặc trưng được chọn ngẫu nhiên thường là $m = \lfloor\sqrt{p}\rfloor$ cho bài toán phân lớp.

Dự đoán cuối cùng được xác định bằng bỏ phiếu đa số (majority voting):
\begin{equation*}
    \hat{y}(\mathbf{x}) = \text{mode}\left\{ h_b(\mathbf{x}) \right\}_{b=1}^{B}
\end{equation*}

Phương sai của mô hình tập hợp được giới hạn bởi:
\begin{equation*}
    \text{Var}\left[\hat{y}\right] = \rho \sigma^2 + \frac{1 - \rho}{B} \sigma^2
\end{equation*}
trong đó $\rho$ là hệ số tương quan trung bình giữa các cây và $\sigma^2$ là phương sai của mỗi cây đơn lẻ. Việc chọn ngẫu nhiên $m \ll p$ đặc trưng tại mỗi nút giúp giảm $\rho$, từ đó giảm phương sai tổng thể.\\

\textbf{Ưu điểm:}
\begin{itemize}
    \item \textbf{Độ bền vững với ngoại lai (Robustness):} Cơ chế trung bình hóa từ hàng trăm cây quyết định độc lập giúp mô hình miễn nhiễm với các giá trị bất thường ngẫu nhiên (random spikes) của cảm biến.
    \item \textbf{Huấn luyện song song:} Các cây hoàn toàn độc lập nên có thể huấn luyện song song trên nhiều lõi CPU, tận dụng hiệu quả kiến trúc đa nhân.
    \item \textbf{Ước lượng lỗi ngoài túi (Out-of-Bag Error):} Khoảng $37\%$ mẫu không xuất hiện trong mỗi bootstrap có thể được dùng làm tập kiểm tra nội bộ, cung cấp ước lượng lỗi tổng quát hóa mà không cần tập validation riêng.
    \item \textbf{Đo lường tầm quan trọng đặc trưng (Feature Importance):} Cung cấp thước đo dựa trên mức giảm tạp chất trung bình (Mean Decrease in Impurity) hoặc hoán vị (Permutation Importance), hỗ trợ giải thích mô hình.
\end{itemize}

\textbf{Nhược điểm:}
\begin{itemize}
    \item \textbf{Độ nhạy thấp với lớp thiểu số:} Thuật toán Bagging hội tụ về phía dự đoán nhãn đa số do mỗi bootstrap chỉ chứa một tỷ lệ nhỏ mẫu lớp thiểu số. Cần kết hợp với kỹ thuật Balanced Random Forest (lấy mẫu phân tầng trong bootstrap) để cải thiện.
    \item \textbf{Kích thước mô hình lớn:} Với $B$ cây sâu không giới hạn, file mô hình có thể lên tới hàng GB trên ổ cứng, gây khó khăn cho việc triển khai và lưu trữ.
    \item \textbf{Không có cơ chế học tuần tự:} Khác với Boosting, mỗi cây không học từ sai sót của cây trước, nên không tập trung sửa lỗi ở vùng dữ liệu khó — điểm yếu quan trọng trong bài toán mất cân bằng.
\end{itemize}

\paragraph{Extra Trees (Extremely Randomized Trees)}\hspace{1cm}

Extra Trees (Extremely Randomized Trees) là một biến thể cực đoan hóa của Random Forest. Khác biệt cốt lõi nằm ở hai điểm: Sử dụng toàn bộ tập huấn luyện $\mathcal{D}$ thay vì lấy mẫu bootstrap và ngẫu nhiên hóa hoàn toàn cả ngưỡng phân chia tại mỗi nút.

Cụ thể, tại mỗi nút, thuật toán chọn ngẫu nhiên $K$ đặc trưng. Với mỗi đặc trưng $j$ được chọn, ngưỡng chia $s_j$ được lấy ngẫu nhiên đồng đều trong khoảng giá trị:
\begin{equation*}
    s_j \sim \mathcal{U}\left(\min_{i} x_i^j, \; \max_{i} x_i^j\right)
\end{equation*}
Sau đó, cặp $(j^*, s^*)$ cho điểm chia tốt nhất được chọn theo tiêu chí Gini hoặc Entropy:
\begin{equation*}
    (j^*, s^*) = \arg\max_{j \in \{1,\ldots,K\}, \; s_j} \; \text{Gain}(j, s_j)
\end{equation*}

Việc ngẫu nhiên hóa ngưỡng chia khiến hệ số tương quan $\rho$ giữa các cây giảm sâu hơn so với Random Forest, dẫn đến phương sai tổng thể thấp hơn theo công thức $\text{Var}[\hat{y}] = \rho \sigma^2 + \frac{1-\rho}{B}\sigma^2$.\\

\textbf{Ưu điểm:}
\begin{itemize}
    \item \textbf{Tốc độ huấn luyện nhanh hơn Random Forest:} Không cần tìm kiếm ngưỡng chia tối ưu cho mỗi đặc trưng (chỉ chọn ngẫu nhiên), giúp giảm đáng kể chi phí tính toán tại mỗi nút từ $O(n \log n)$ xuống $O(n)$.
    \item \textbf{Giảm phương sai triệt để hơn:} Hệ số tương quan $\rho$ giữa các cây thấp hơn Random Forest nhờ mức độ ngẫu nhiên hóa cao hơn, giúp mô hình tập hợp có phương sai tổng thể thấp hơn.
    \item \textbf{Kháng overfitting mạnh:} Ngưỡng chia ngẫu nhiên ngăn cây học quá sát các mẫu nhiễu cục bộ. Sử dụng toàn bộ dữ liệu (không bootstrap) giúp mỗi cây có cái nhìn đầy đủ hơn về phân phối dữ liệu, giảm thiên kiến (bias).
    \item \textbf{Phù hợp với dữ liệu nhiễu cao:} Đặc tính ngẫu nhiên hóa cực đoan giúp Extra Trees miễn nhiễm tốt với nhiễu cảm biến, phù hợp với bản chất dữ liệu đo lường năng lượng.
\end{itemize}

\textbf{Nhược điểm:}
\begin{itemize}
    \item \textbf{Thiên kiến (Bias) cao hơn Random Forest:} Ngưỡng chia ngẫu nhiên thay vì tối ưu có thể tạo ra các phân vùng kém chính xác hơn, khiến mỗi cây đơn lẻ có sai số thiên kiến cao hơn.
    \item \textbf{Nhạy cảm với đặc trưng không liên quan:} Do không tìm kiếm ngưỡng tối ưu, thuật toán khó phân biệt giữa các đặc trưng hữu ích và đặc trưng nhiễu, đặc biệt khi tập đặc trưng có chiều cao.
    \item \textbf{Cùng hạn chế Bagging với lớp thiểu số:} Tương tự Random Forest, Extra Trees không có cơ chế tập trung vào các mẫu khó (hard samples) của lớp thiểu số, kém hiệu quả hơn Boosting trong bài toán mất cân bằng nghiêm trọng.
\end{itemize}

\subsubsection{Nhóm Phương Pháp Không Giám Sát (Unsupervised) }
Được ứng dụng để nhận diện bất thường dựa trên cấu trúc tô-pô và phân phối thống kê của dữ liệu, không phụ thuộc vào nhãn. Trong nhóm thuật toán này, chỉ có thuật toán Isolation Forest là khả thi để thực hiện.\\

\textbf{Isolation Forest}

Phương pháp tiếp cận dựa trên sự cô lập. Điểm dữ liệu càng dễ bị tách biệt khỏi quần thể chung bằng các mặt phẳng ngẫu nhiên thì điểm số bất thường càng cao. Mỗi Isolation Tree (iTree) thực hiện phân chia ngẫu nhiên đệ quy: chọn ngẫu nhiên một đặc trưng $q$ và một ngưỡng $p$ trong khoảng $[\min(x_q), \max(x_q)]$.

Điểm số bất thường (Anomaly Score) của điểm $\mathbf{x}$ được xác định bởi:
\begin{equation*}
    s(\mathbf{x}, n) = 2^{-\frac{E[h(\mathbf{x})]}{c(n)}}
\end{equation*}
trong đó $E[h(\mathbf{x})]$ là chiều dài đường đi trung bình (average path length) của $\mathbf{x}$ trên tập hợp các iTree và $c(n)$ là chiều dài đường đi trung bình của phép tìm kiếm không thành công trong cây nhị phân BST với $n$ nút:
\begin{equation*}
    c(n) = 2H(n-1) - \frac{2(n-1)}{n}
\end{equation*}
với $H(i) = \ln(i) + \gamma_E$ là số điều hòa (Harmonic number) và $\gamma_E \approx 0.5772$ là hằng số Euler-Mascheroni. Khi $s \to 1$, điểm dữ liệu có khả năng cao là bất thường; khi $s \to 0.5$, điểm dữ liệu thuộc quần thể bình thường.\\

\textbf{Ưu điểm:}
\begin{itemize}
    \item \textbf{Độ phức tạp tuyến tính:} Độ phức tạp thời gian đạt $O(t \cdot \psi \cdot \log \psi)$ cho việc xây dựng $t$ cây với kích thước không gian mẫu phụ $\psi$ và $O(t \cdot \log \psi)$ cho việc tính điểm bất thường mỗi mẫu. Thuật toán mở rộng cực kỳ hiệu quả trên dữ liệu khổng lồ.
    \item \textbf{Không cần giả định phân phối:} Khác với các phương pháp dựa trên mật độ (LOF) hay phân phối (GMM), Isolation Forest không yêu cầu bất kỳ giả định nào về phân phối dữ liệu, hoạt động tốt trên dữ liệu có phân phối phức tạp.
    \item \textbf{Khả năng mở rộng (Scalability):} Tham số $\psi$ (sub-sample size) cho phép kiểm soát trực tiếp lượng dữ liệu mỗi cây sử dụng, giữ bộ nhớ và thời gian trong tầm kiểm soát ngay cả với hàng chục triệu dòng dữ liệu.
\end{itemize}

\textbf{Nhược điểm:}
\begin{itemize}
    \item \textbf{Hạn chế về bất thường ngữ cảnh (Contextual Anomalies):} Chỉ hiệu quả trong việc tìm kiếm bất thường toàn cục (Global outliers). Có thể bỏ qua các bất thường mang tính chu kỳ của dữ liệu chuỗi thời gian nếu không có kỹ thuật trích xuất đặc trưng độ trễ (Lag features) phù hợp.
    \item \textbf{Thiên kiến với đặc trưng có nhiều giá trị phân biệt:} Phép chia ngẫu nhiên có xu hướng ưu tiên các đặc trưng có phạm vi giá trị rộng, có thể dẫn đến điểm bất thường bị sai lệch nếu không chuẩn hóa dữ liệu.
    \item \textbf{Không có cơ chế học có giám sát:} Anomaly Score được tính hoàn toàn dựa trên cấu trúc dữ liệu, không tận dụng được nhãn đã có. Do đó trong pipeline, Isolation Forest chỉ đóng vai trò tạo meta-feature chứ không phải mô hình phân lớp chính.
\end{itemize}

\subsubsection{Kết luận và Thiết kế Pipeline}
Dựa trên phân tích ưu nhược điểm của từng thuật toán, nhóm nghiên cứu đề xuất một pipeline gồm bốn giai đoạn, tuân theo nguyên tắc: huấn luyện rời rạc từng mô hình trước, tối ưu hóa riêng biệt, sau đó kết hợp (ensemble) để đạt hiệu suất tối đa.

\begin{enumerate}
    \item \textbf{Giai đoạn 1: Tiền xử lý và Trích xuất đặc trưng (Feature Engineering)}
    \begin{itemize}
        \item Thực hiện các bước làm sạch dữ liệu: xử lý giá trị thiếu, loại bỏ ngoại lai rõ ràng, chuẩn hóa các biến liên tục.
        \item Trích xuất các đặc trưng chuỗi thời gian: cửa sổ trượt (rolling window), đặc trưng độ trễ (lag features), mã hóa chu kỳ (cyclical encoding) cho giờ/ngày/tháng.
        \item Ứng dụng \textbf{Isolation Forest} nhằm quét toàn bộ không gian dữ liệu. Điểm số bất thường $s(\mathbf{x}, n)$ được ghép nối vào tập dữ liệu gốc đóng vai trò như một đặc trưng siêu dữ liệu (Meta-feature), cung cấp thêm tín hiệu hội tụ cho các thuật toán phân lớp ở các giai đoạn sau.
    \end{itemize}
    
    \item \textbf{Giai đoạn 2: Huấn luyện rời rạc từng mô hình (Baseline Modeling)}
    \begin{itemize}
        \item Tiến hành huấn luyện độc lập từng thuật toán với các siêu tham số mặc định hoặc hợp lý ban đầu. Các mô hình được đánh giá bao gồm:
        \begin{enumerate}
            \item \textbf{LightGBM} — tận dụng tốc độ và hiệu quả bộ nhớ trên dữ liệu quy mô lớn.
            \item \textbf{XGBoost} — khai thác cơ chế điều chuẩn $L_1/L_2$ mạnh mẽ.
            \item \textbf{CatBoost} — khai thác khả năng xử lý biến định danh và cấu trúc Ordered Boosting.
            \item \textbf{Random Forest} — đánh giá sức mạnh của phương pháp Bagging với cơ chế chọn đặc trưng ngẫu nhiên.
            \item \textbf{Extra Trees} — đánh giá hiệu quả của ngẫu nhiên hóa cực đoan trong việc giảm phương sai.
        \end{enumerate}
        \item Đánh giá hiệu năng từng mô hình đơn lẻ dựa trên các độ đo phù hợp với dữ liệu mất cân bằng: \textbf{F1-Score}, \textbf{PR-AUC} (Precision-Recall AUC) và \textbf{ROC-AUC}. Giai đoạn này thiết lập mức hiệu suất cơ sở (baseline) cho từng thuật toán.
    \end{itemize}

    \item \textbf{Giai đoạn 3: Tối ưu hóa từng mô hình (Hyperparameter Tuning)}
    \begin{itemize}
        \item \textbf{Tinh chỉnh siêu tham số:} Sử dụng phương pháp tìm kiếm Bayesian (Optuna) hoặc Random Search kết hợp với Cross-Validation phân tầng (Stratified K-Fold) để tìm bộ siêu tham số tối ưu cho từng mô hình riêng biệt.
        \item \textbf{Tối ưu hóa ngưỡng cắt xác suất (Threshold Moving):} Thay vì sử dụng ngưỡng mặc định $0.5$, thực hiện tìm kiếm ngưỡng tối ưu trên tập validation nhằm tối đa hóa F1-Score, thích ứng với sự khan hiếm cực độ của nhãn $1$.
        \item \textbf{Xử lý mất cân bằng lớp:} Áp dụng các kỹ thuật đặc thù cho từng thuật toán — \texttt{scale\_pos\_weight} (XGBoost), \texttt{is\_unbalance} (LightGBM), \texttt{auto\_class\_weights} (CatBoost), \texttt{class\_weight='balanced'} (Random Forest, Extra Trees).
        \item \textbf{So sánh kết quả:} Đối chiếu hiệu năng của từng mô hình sau tối ưu hóa so với baseline ở Giai đoạn 2, xác định top mô hình đơn lẻ tốt nhất.
    \end{itemize}

    \item \textbf{Giai đoạn 4: Dung hợp mô hình — Lắp ráp Ensemble (Model Ensemble)}
    \begin{itemize}
        \item \textbf{Stacking (Xếp chồng):} Sử dụng dự đoán xác suất (probability predictions) từ các mô hình đã tối ưu ở Giai đoạn 3 làm đặc trưng đầu vào cho một mô hình siêu học (meta-learner). Meta-learner học cách kết hợp tối ưu các tín hiệu từ nhiều kiến trúc khác nhau:
        \begin{equation*}
            \hat{y}_{\text{stack}} = g\left(\hat{p}_{\text{LGBM}}, \hat{p}_{\text{XGB}}, \hat{p}_{\text{Cat}}, \hat{p}_{\text{RF}}, \hat{p}_{\text{ET}}\right)
        \end{equation*}
        trong đó $g(\cdot)$ là meta-learner (thường là Logistic Regression hoặc một LightGBM nhẹ), $\hat{p}$ là xác suất dự đoán từ mỗi mô hình cơ sở.
        \item \textbf{Blending (Trộn):} Kết hợp dự đoán thông qua trung bình có trọng số:
        \begin{equation*}
            \hat{p}_{\text{blend}} = \sum_{m=1}^{M} w_m \cdot \hat{p}_m, \quad \text{với} \quad \sum_{m=1}^{M} w_m = 1
        \end{equation*}
        trong đó trọng số $w_m$ được tối ưu hóa trên tập validation để tối đa hóa F1-Score.
        \item \textbf{Soft Voting:} Tính trung bình cộng xác suất từ tất cả các mô hình, sau đó áp dụng ngưỡng cắt đã tối ưu:
        \begin{equation*}
            \hat{y}_{\text{vote}} = \mathbb{1}\left[\frac{1}{M}\sum_{m=1}^{M} \hat{p}_m \geq \tau^*\right]
        \end{equation*}
        trong đó $\tau^*$ là ngưỡng tối ưu được xác định từ Giai đoạn 3.
        \item \textbf{Mixture of Experts theo ngữ cảnh:} Thay vì gán một bộ trọng số cố định cho mọi mẫu, huấn luyện một bộ định tuyến (gating model) để dự đoán trọng số của từng mô hình cơ sở dựa trên đặc trưng ngữ cảnh như \texttt{building\_id}, \texttt{meter}, thời điểm và điểm bất thường từ Isolation Forest. Xác suất dự đoán kết hợp được tính như sau:
        \begin{equation*}
            \hat{p}_{\text{MoE}}(\mathbf{x}) = \sum_{m=1}^{M} g_m(\mathbf{z})\,\hat{p}_m(\mathbf{x}), \quad \sum_{m=1}^{M} g_m(\mathbf{z}) = 1
        \end{equation*}
        trong đó $\mathbf{z}$ là các đặc trưng ngữ cảnh và $g_m(\mathbf{z})$ là trọng số do gating model sinh ra cho chuyên gia thứ $m$. Ví dụ, bộ định tuyến có thể ưu tiên CatBoost ở nhóm công tơ có nhiều biến phân loại, đồng thời ưu tiên LightGBM ở các nhóm dữ liệu lớn.
        \item \textbf{Huấn luyện chống rò rỉ:} Tạo dự đoán out-of-fold (OOF) cho các mô hình cơ sở trên tập huấn luyện để huấn luyện gating model; không dùng dự đoán in-sample vì chúng có thể khiến bộ định tuyến học thuộc lỗi huấn luyện. Với dữ liệu chuỗi thời gian, các fold phải được chia theo thứ tự thời gian (hoặc theo nhóm thiết bị và thời gian), không xáo trộn ngẫu nhiên, để tránh thông tin tương lai lọt vào quá trình huấn luyện.
        \item \textbf{So sánh tổng thể:} Đối chiếu hiệu năng của các phương pháp Ensemble (Stacking, Blending, Soft Voting và Mixture of Experts) so với mô hình đơn lẻ tốt nhất ở Giai đoạn 3. Báo cáo các độ đo F1-Score, PR-AUC, ROC-AUC cùng thời gian suy luận để đánh giá đồng thời chất lượng dự đoán và chi phí triển khai.
    \end{itemize}
\end{enumerate}



\subsubsection{Luận giải việc loại trừ các thuật toán không được chọn vào Pipeline}
Trong quá trình xây dựng quy trình triển khai (pipeline), mặc dù nhiều thuật toán đã được xem xét ở khâu khảo sát lý thuyết, nhưng phần lớn đã bị loại bỏ ở khâu thực nghiệm do vấp phải những rào cản chí mạng về tính toán và bản chất dữ liệu. Dưới đây là phân tích nguyên nhân loại trừ đối với từng nhóm thuật toán cụ thể:

\textbf{1. Nhóm không khả thi về Độ phức tạp tính toán (Computational Bottleneck)}
\begin{itemize}
    \item \textbf{Các thuật toán bị loại:} One-Class SVM, Local Outlier Factor (LOF) và DBSCAN.
    \item \textbf{Lý do loại trừ:} Cả ba thuật toán này đều đòi hỏi tính toán ma trận khoảng cách hoặc tối ưu hóa biên không gian với độ phức tạp thời gian và bộ nhớ xấp xỉ $\mathcal{O}(n^2)$. Với quy mô dữ liệu của bài toán lên tới hàng chục triệu dòng (hơn 20 triệu quan trắc), việc chạy các thuật toán này trên thiết lập phần cứng thông thường (thậm chí trên Kaggle Notebooks với 30GB RAM) sẽ ngay lập tức gây ra lỗi tràn bộ nhớ (Out-Of-Memory) hoặc thời gian huấn luyện kéo dài nhiều tuần lễ mà không thể hội tụ.
\end{itemize}

\textbf{2. Nhóm nhạy cảm cực độ với nhiễu dữ liệu (Noise Sensitivity)}
\begin{itemize}
    \item \textbf{Thuật toán bị loại:} AdaBoost.
    \item \textbf{Lý do loại trừ:} Dữ liệu đo lường năng lượng từ các cảm biến luôn đi kèm với tỷ lệ nhiễu ngẫu nhiên (sensor noise/glitches) rất cao. Cơ chế của AdaBoost là phạt (penalty) các mẫu bị đoán sai bằng cách tăng trọng số theo cấp số nhân. Điều này khiến thuật toán bị đánh lừa: thay vì học cách phát hiện bất thường thực sự, nó lại tập trung khớp (fit) các điểm nhiễu ngẫu nhiên, dẫn đến hiện tượng quá khớp (overfitting) nghiêm trọng và giảm khả năng tổng quát hóa trên tập kiểm tra.
\end{itemize}

\textbf{3. Nhóm giới hạn về Không gian phân phối và Giả định tuyến tính}
\begin{itemize}
    \item \textbf{Các thuật toán bị loại:} K-Means, Gaussian Mixture Models (GMM) và PCA.
    \item \textbf{Lý do loại trừ:} 
    \begin{itemize}
        \item \textbf{K-Means \& GMM:} Gặp khó khăn lớn khi chiều dữ liệu tăng lên (do trích xuất thêm các đặc trưng trễ - lag features). K-Means bị bó buộc bởi giả định các cụm hình cầu, trong khi GMM rất dễ kẹt ở cực tiểu địa phương. Cả hai đều thiên về tìm kiếm bất thường toàn cục mà bỏ sót các \textit{bất thường ngữ cảnh} (Contextual anomalies) thường thấy trong chuỗi thời gian tiêu thụ điện.
        \item \textbf{PCA:} Phương pháp này chỉ mô phỏng được các mối quan hệ tuyến tính bằng cách chiếu dữ liệu lên các thành phần trực giao. Tuy nhiên, tương quan giữa tiêu thụ năng lượng, nhiệt độ thời tiết và thói quen sinh hoạt là một hàm phi tuyến tính phức tạp, khiến sai số tái tạo của PCA không phản ánh đúng bản chất của sự bất thường.
    \end{itemize}
\end{itemize}

\textbf{Tiểu kết:} Việc loại trừ các thuật toán trên không phủ nhận giá trị của chúng trong lĩnh vực Học máy, mà là kết quả tất yếu của quá trình chọn lọc tự nhiên để phù hợp với ba đặc tính cốt lõi của bài toán này: Tính siêu lớn của dữ liệu, sự mất cân bằng nhãn cực đoan và bản chất phi tuyến tính của chuỗi thời gian năng lượng.


\subsubsection{Các kỹ thuật và thuật toán bổ trợ khác (Có thể tích hợp thêm)}
Để tối ưu hóa sâu hơn cho mô hình và dữ liệu, một số kỹ thuật và thuật toán bổ trợ sau có thể được kết hợp vào pipeline:
\begin{itemize}
    \item \textbf{SMOTE (Synthetic Minority Over-sampling Technique):} Kỹ thuật sinh mẫu nhân tạo nhằm giải quyết mất cân bằng lớp ở mức độ dữ liệu (Data-level). Đối với tập dữ liệu hàng chục triệu dòng, có thể áp dụng trên một tập mẫu con (sub-sample) hoặc kết hợp với Undersampling (như SMOTE-ENN) để cải thiện độ phân tách mà không gây tràn bộ nhớ.
    \item \textbf{Optuna (Hyperparameter Optimization Framework):} Framework tối ưu hóa siêu tham số tự động dựa trên Bayesian Optimization (cụ thể là Tree-structured Parzen Estimator - TPE). Optuna hỗ trợ cắt tỉa (pruning) các thử nghiệm không hứa hẹn, giúp tiết kiệm tối đa thời gian tính toán khi huấn luyện các mô hình Boosting lớn.
\end{itemize}

\subsubsection{Chi tiết Pipeline Tiền xử lý dữ liệu và Áp dụng trên File}
Dựa trên cấu trúc bộ dữ liệu, quy trình tiền xử lý sẽ tác động trực tiếp lên các tệp (nằm trong thư mục \texttt{data/}) theo từng bước cụ thể sau:

\textbf{1. Tiền xử lý tập dữ liệu đặc trưng (Feature Processing):}
\begin{itemize}
    \item \textbf{Files áp dụng:} \texttt{train\_features.csv} (dành cho huấn luyện) và \texttt{test\_features.csv} (dành cho dự đoán/đánh giá). Vì hai file này đã chứa đầy đủ dữ liệu đo lường (\texttt{meter\_reading}) và nhãn (\texttt{anomaly}), ta không cần phải ghép nối với các file gốc.
    \item \textbf{Xử lý nhiễu (Noise):} Loại bỏ hoặc sửa chữa các giá trị \texttt{meter\_reading} bất thường rõ ràng do lỗi cảm biến (ví dụ: các chuỗi giá trị 0 liên tiếp bất thường ở một số tòa nhà).
    \item \textbf{Xử lý giá trị thiếu (Missing Values):} Điền khuyết các biến thời tiết (\texttt{air\_temperature}, \texttt{dew\_temperature}, v.v.) bằng nội suy tuyến tính (Linear Interpolation) hoặc giá trị trung bình theo nhóm (group by \texttt{site\_id}).
    \item \textbf{Xử lý biến định danh (Categorical Variables):} Các biến \texttt{primary\_use}, \texttt{building\_id}, \texttt{site\_id} sẽ được mã hóa bằng Target Encoding (được tích hợp sẵn trong CatBoost/LightGBM) hoặc Label Encoding.
    \item \textbf{Đặc trưng chuỗi thời gian (Time-series Features):} Tạo các đặc trưng về chu kỳ (Cyclical Features) từ cột \texttt{timestamp} (như $\sin/\cos$ của \texttt{hour}, \texttt{month}) và khai thác lại các đặc trưng độ trễ (lag features) có sẵn (ví dụ: \texttt{air\_temperature\_mean\_lag7}).
\end{itemize}

\textbf{2. Trích xuất siêu đặc trưng bất thường (Anomaly Meta-feature):}
\begin{itemize}
    \item Ứng dụng \textbf{Isolation Forest} trên tập huấn luyện (sau khi xử lý nhiễu và điền khuyết) để sinh ra \texttt{anomaly\_score}. Đặc trưng mới này sau đó được bổ sung như một biến đầu vào (input feature) cho mô hình phân lớp ở giai đoạn tiếp theo.
\end{itemize}

\subsubsection{Tài liệu tham khảo Quốc tế}
Danh sách các bài báo khoa học chuẩn mực đã được tham khảo làm nền tảng lý thuyết cho việc lựa chọn thuật toán và thiết kế pipeline:
\begin{itemize}
    \item \textbf{XGBoost:} Chen, T., \& Guestrin, C. (2016). XGBoost: A Scalable Tree Boosting System. \textit{KDD '16}. DOI: \url{https://doi.org/10.1145/2939672.2939785}
    \item \textbf{LightGBM:} Ke, G., Meng, Q., Finley, T., et al. (2017). LightGBM: A Highly Efficient Gradient Boosting Decision Tree. \textit{NeurIPS 2017}. URL: \url{https://papers.nips.cc/paper/6907-lightgbm-a-highly-efficient-gradient-boosting-decision-tree.pdf}
    \item \textbf{CatBoost:} Prokhorenkova, L., Gusev, G., Vorobev, A., et al. (2018). CatBoost: unbiased boosting with categorical features. \textit{NeurIPS 2018}. DOI: \url{https://doi.org/10.5555/3295222.3295328}
    \item \textbf{Random Forest:} Breiman, L. (2001). Random Forests. \textit{Machine Learning}, 45(1), 5–32. DOI: \url{https://doi.org/10.1023/A:1010933404324}
    \item \textbf{Extra Trees:} Geurts, P., Ernst, D., \& Wehenkel, L. (2006). Extremely randomized trees. \textit{Machine Learning}, 63(1), 3–42. DOI: \url{https://doi.org/10.1007/s10994-006-6226-1}
    \item \textbf{Isolation Forest:} Liu, F. T., Ting, K. M., \& Zhou, Z.-H. (2008). Isolation Forest. \textit{ICDM 2008}. DOI: \url{https://doi.org/10.1109/icdm.2008.17}
    \item \textbf{SMOTE:} Chawla, N. V., et al. (2002). SMOTE: Synthetic Minority Over-sampling Technique. \textit{JAIR}, 16, 321–357. DOI: \url{https://doi.org/10.1613/jair.953}
    \item \textbf{Optuna:} Akiba, T., et al. (2019). Optuna: A Next-generation Hyperparameter Optimization Framework. \textit{KDD '19}. DOI: \url{https://doi.org/10.1145/3292500.3330701}
    \item \textbf{Dữ liệu ASHRAE:} Miller, C., et al. (2020). The ASHRAE Great Energy Predictor III competition: Overview and results. \textit{Science and Technology for the Built Environment}. DOI: \url{https://doi.org/10.1080/23744731.2020.1795514}
\end{itemize}

\subsection{Các phương pháp Học sâu hiện có để giải quyết bài toán}
%Ghi rõ và ưu nhược điểm